\begin{frame}{CV 추정값은 ``진짜 성능''이 아니다}
  \begin{block}{CV 추정값(12.75)이 정직한 값(8.25)보다 53\%나 높다}
    각 fold의 ``train 평균''(6.5, 6.0, 5.5, 5.0, 4.5)이 진짜 평균 5.5
    주위를 흔들리기 때문. train 평균이 1.0만큼 어긋난 fold 0에서 그
    어긋남이 val 오차에 \emph{그대로} 더쳐진다.
    \vskip0.3em
    \textbf{CV 추정값은 ``진짜 성능''이 아니라 ``검증 절차의 시뮬레이션''} —
    절차 안의 작은 흔들림도 포함되므로 특정 모델 하나엔 오히려
    \emph{높게}(비관적으로) 나올 수도.
  \end{block}
  \vskip0.4em
  \begin{itemize}
    \item ``CV 수치가 항상 낙관적''이라는 직관을 \textbf{바로잡}자 —
      \kq{CV의 낙관적 편향은 ``한 모델을 평가''할 때 생기는 게 아니라,
      ``여러 후보 중 \emph{제일 좋은} 모델을 고르는'' 다음 단계에서
      생긴다}(§최적화 편향).
    \item 모델 \emph{하나}를 단순히 평가하는 단계(여기처럼 유연하지
      않으면) — CV가 진짜 오차보다 \emph{낮게} 나올 리가 없고, 실제로
      \emph{높게} 나올 수도.
  \end{itemize}
\end{frame}
